%==============================================================================
%  CHAPITRE 1 — INTRODUCTION GÉNÉRALE
%==============================================================================
\chapter{Introduction Générale}
\label{chap:introduction}
\minitoc

%==============================================================================
\section{Contexte et motivation}
\label{sec:contexte}

L'enseignement à distance a connu, au cours de la dernière décennie, une
transformation profonde portée par l'essor des plateformes d'apprentissage en
ligne ouvertes et massives, communément désignées par l'acronyme anglais MOOC
(\emph{Massive Open Online Courses}). Depuis le lancement des premiers cours à
large échelle au début des années 2010, ces dispositifs ont permis de
démocratiser l'accès à la connaissance en offrant à des millions d'apprenants,
sans condition géographique ni, le plus souvent, financière, des contenus
pédagogiques produits par des établissements de premier plan
\cite{breslow2013studying,daniel2012making}. La crise sanitaire de 2020 a encore
accéléré cette dynamique, faisant de l'apprentissage en ligne une composante
désormais structurelle des systèmes éducatifs contemporains.

Cette promesse de massification se heurte toutefois à un constat largement
documenté : les MOOC souffrent de taux d'abandon particulièrement élevés. La
proportion d'apprenants menant un cours à son terme dépasse rarement dix à
quinze pour cent des inscrits, et la majorité des décrochages survient dès les
premières semaines \cite{onah2014dropout,jordan2015massive}. Ce phénomène,
souvent qualifié de paradoxe des MOOC, oppose une accessibilité sans précédent à
une efficacité pédagogique réelle bien plus modeste qu'attendu. Comprendre et
réduire ce décrochage constitue, depuis lors, l'un des enjeux majeurs de la
recherche en technologies éducatives.

Parmi les multiples facteurs explicatifs de l'abandon — manque de temps,
prérequis insuffisants, absence de reconnaissance institutionnelle — un
déterminant occupe une place centrale : la baisse de l'\textbf{engagement} de
l'apprenant \cite{kizilcec2013deconstructing,hew2016promoting}. L'engagement
recouvre la mobilisation cognitive, comportementale et émotionnelle de
l'individu vis-à-vis de la tâche d'apprentissage. Or, dans une salle de classe
traditionnelle, un enseignant perçoit en permanence des signaux fins — un regard
fuyant, une posture affaissée, une expression de perplexité — et adapte
spontanément son discours : il reformule, ralentit, propose un exemple, relance
l'attention. Cette boucle de régulation, fondamentale pour l'apprentissage,
disparaît presque entièrement dans les environnements en ligne de masse, où
l'apprenant fait face seul à un contenu figé, identique pour tous et insensible
à son état du moment.

Cette absence de rétroaction affective n'est pas anodine. Les travaux en
psychologie de l'apprentissage et en informatique affective
(\emph{affective computing}) ont établi que les émotions ne sont pas
périphériques à la cognition mais en constituent un rouage essentiel
\cite{picard1997affective,pekrun2006control}. Des états tels que l'ennui, la
confusion ou la frustration influencent directement la profondeur du traitement
de l'information, la persévérance et la mémorisation
\cite{dmello2012dynamics,baker2010better}. Détecter ces états et y répondre de
manière appropriée apparaît dès lors comme une voie prometteuse pour restaurer,
au moins partiellement, la régulation pédagogique perdue dans les MOOC.

Deux évolutions technologiques rendent aujourd'hui cette ambition crédible. La
première est la maturité des méthodes de \textbf{détection automatique de
l'engagement} à partir de signaux multimodaux. L'analyse des expressions
faciales, de la posture, de la voix ou de signaux physiologiques permet
désormais d'inférer, avec une fiabilité croissante, l'état affectif d'un
apprenant \cite{dmello2015review,whitehill2014faces}. La diffusion de corpus
annotés de référence, au premier rang desquels \textsc{DAiSEE}
\cite{gupta2016daisee} pour l'engagement et \textsc{IEMOCAP}
\cite{busso2008iemocap} pour l'émotion vocale, a fourni à la communauté les
ressources nécessaires à l'entraînement et à l'évaluation de tels modèles.

La seconde évolution est l'avènement de l'\textbf{intelligence artificielle
générative}, et plus particulièrement des grands modèles de langue
(\emph{Large Language Models}, LLM). Ces modèles sont capables de produire, à la
demande et en langage naturel, des explications, des reformulations, des
analogies ou des encouragements de qualité, ouvrant la perspective d'une
génération de contenu pédagogique personnalisée à grande échelle
\cite{brown2020language,kasneci2023chatgpt}. Là où la personnalisation
supposait jusqu'ici de produire et de maintenir manuellement d'innombrables
variantes de contenu, un LLM correctement guidé peut générer, en quelques
secondes, une réponse adaptée à un état d'apprentissage particulier.

C'est à la croisée de ces deux courants que se situe le présent travail. L'idée
directrice consiste à réunir, au sein d'une même chaîne de traitement, la
\emph{perception} de l'engagement émotionnel et la \emph{génération} adaptative
de contenu, afin d'esquisser une forme de régulation pédagogique automatique
inspirée de celle de l'enseignant humain. C'est cette intuition que matérialise
le prototype \emolearn{}, objet de ce mémoire.

%==============================================================================
\section{Problématique}
\label{sec:problematique}

Le constat précédent fait apparaître un écart à combler. D'un côté, les MOOC
disposent d'une audience massive mais d'une capacité d'adaptation quasi nulle aux
états affectifs de leurs apprenants. De l'autre, les briques technologiques —
détection multimodale de l'engagement, génération de contenu par LLM — existent
mais demeurent largement étudiées de manière isolée. La question centrale qui
structure ce mémoire peut dès lors s'énoncer ainsi :

\begin{quote}
\itshape
Comment concevoir un dispositif capable d'inférer automatiquement l'état
d'engagement d'un apprenant à partir de signaux multimodaux, puis de générer en
réponse un contenu pédagogique adapté à cet état, tout en garantissant la
cohérence et la traçabilité de la décision pédagogique sous-jacente ?
\end{quote}

Cette interrogation générale se décline en plusieurs questions de recherche plus
précises, qui guideront l'ensemble du travail :

\begin{description}
  \item[QR1 — Perception.] Dans quelle mesure des signaux multimodaux
  hétérogènes (expressions faciales issues de \textsc{DAiSEE}, signal vocal issu
  de \textsc{IEMOCAP}) peuvent-ils être combinés par une stratégie de fusion pour
  estimer l'état d'engagement de manière plus fiable qu'une modalité unique, et
  quel est le coût de cette fusion en termes de compromis entre exactitude
  globale et sensibilité aux états critiques ? Dans quelle mesure, en outre,
  l'apprentissage supervisé de la correspondance entre les sorties de modèles
  pré-entraînés et la taxonomie d'engagement visée améliore-t-il cette perception
  par rapport à une mise en correspondance définie \emph{a priori} ?

  \item[QR2 — Décision.] Comment formaliser une correspondance pédagogiquement
  justifiable entre les états d'engagement détectés et les stratégies
  d'intervention, c'est-à-dire un modèle \emph{engagement\,$\rightarrow$\,adaptation}
  reliant chaque état à un type de contenu et à un patron de génération précis ?

  \item[QR3 — Génération.] Dans quelle mesure un grand modèle de langue, guidé
  par une bibliothèque de patrons de prompts structurés, peut-il produire un
  contenu pédagogique pertinent, clair et effectivement adapté à l'état
  d'engagement visé ?

  \item[QR4 — Validation.] Selon quel protocole peut-on évaluer, de façon
  honnête et multi-niveaux, à la fois la qualité de la perception, la pertinence
  du modèle de décision et la qualité du contenu généré, compte tenu des
  contraintes propres à un travail de master ?
\end{description}

Au-delà de ces questions, le travail assume d'emblée plusieurs
\textbf{tensions conceptuelles} qui en délimitent le périmètre et qu'il convient
d'expliciter dès l'introduction, par souci de rigueur. Premièrement, \emolearn{}
est un \emph{prototype de simulation} et non un système déployé en temps réel sur
de véritables MOOC : l'engagement y provient de corpus annotés a posteriori et
non d'une captation \emph{live} d'apprenants. Deuxièmement, la dimension
multimodale repose, dans la chaîne de production, sur un appariement entre des
signaux issus de populations distinctes, le visage et la voix ne provenant pas du
même individu ; cette limite est levée uniquement pour la phase de validation, à
l'aide d'un sous-corpus audiovisuel aligné. Troisièmement, la mise en
correspondance d'émotions conversationnelles avec des états d'engagement constitue
une hypothèse de construit, et non une équivalence démontrée. Plutôt que de
masquer ces tensions, ce mémoire fait le choix de les identifier, de les cadrer et,
chaque fois que possible, de les \emph{mesurer} : c'est précisément cette honnêteté
méthodologique qui constitue la valeur scientifique du travail.

%==============================================================================
\section{Objectifs et contributions visées}
\label{sec:objectifs}

\subsection*{Objectif général}

L'objectif général de ce mémoire est de \textbf{concevoir, implémenter et
évaluer un framework de bout en bout}, baptisé \emolearn{}, qui détecte l'état
d'engagement émotionnel d'un apprenant à partir de signaux multimodaux, décide
s'il convient d'intervenir et selon quelle stratégie, puis génère un contenu
pédagogique adaptatif au moyen d'un grand modèle de langue.

\subsection*{Objectifs spécifiques}

Cet objectif général se décompose en objectifs spécifiques (OS) qui structurent
également la progression technique du projet :

\begin{description}
  \item[OS1 — Analyse des données.] Conduire une analyse exploratoire du corpus
  \textsc{DAiSEE} (engagement facial) et du corpus \textsc{IEMOCAP} (émotion
  vocale), caractériser la distribution des états d'engagement et constituer des
  scénarios d'apprentissage représentatifs.

  \item[OS2 — Modèle de décision.] Définir une taxonomie des états d'engagement,
  établir une matrice \emph{engagement\,$\rightarrow$\,besoin pédagogique} et
  formaliser les règles de déclenchement et de sélection de stratégie.

  \item[OS3 — Génération adaptative.] Concevoir une bibliothèque de patrons de
  prompts structurés et un pipeline de génération de contenu par LLM,
  accompagnés de mécanismes de contrôle de qualité.

  \item[OS4 — Fusion multimodale et apprentissage léger.] Implémenter et évaluer
  un moteur de fusion combinant visage et voix ; apprendre, pour chaque modalité,
  la correspondance entre les sorties des modèles de perception pré-entraînés et
  la taxonomie d'engagement, plutôt que de la postuler ; mesurer honnêtement
  l'apport du multimodal par rapport aux modalités unimodales.

  \item[OS5 — Prototype de simulation.] Développer un prototype fonctionnel
  (interface de visualisation et lecteur de parcours d'apprentissage) permettant
  de simuler des parcours adaptatifs et de visualiser les adaptations produites.

  \item[OS6 — Validation multi-niveaux.] Définir et instrumenter un protocole
  d'évaluation couvrant la cohérence du modèle de décision, la qualité du contenu
  généré et l'acceptabilité du dispositif.
\end{description}

\subsection*{Contributions visées}

Les contributions revendiquées par ce travail sont les suivantes :

\begin{enumerate}
  \item Une \textbf{architecture multimodale de bout en bout} cohérente, reliant
  perception, décision et génération au sein d'une même chaîne, et reposant sur
  un moteur de décision et un moteur de fusion uniques, partagés entre la
  production de contenu et le banc d'essai de validation.

  \item Un \textbf{modèle engagement\,$\rightarrow$\,adaptation} explicite,
  associant chaque état d'engagement à un type de contenu et à un patron de
  génération parmi une bibliothèque de vingt patrons structurés, assorti de
  règles temporelles tenant compte de la dynamique d'une séquence.

  \item Un \textbf{apprentissage léger des correspondances} entre modèles
  pré-entraînés et taxonomie d'engagement, conduit sur les deux modalités : un
  \emph{supervised classifier} (\textbf{Random Forest} + \textbf{GridSearchCV}) reliant
  les sorties du module facial aux états \textsc{DAiSEE}, et le fine-tuning de la
  \emph{classification head} d'un modèle \textbf{wav2vec2} afin d'y faire émerger
  la classe \emph{Confusion}, absente du modèle d'origine. Chaque modalité conserve,
  en regard, une correspondance heuristique explicite servant de référence comparative.

  \item Une \textbf{évaluation honnête et instrumentée} de la fusion multimodale,
  qui ne dissimule pas les compromis observés mais les quantifie — notamment
  l'arbitrage entre l'exactitude globale et le rappel des états critiques.

  \item Une \textbf{analyse de sensibilité} des deux hypothèses les plus
  discutables du dispositif — la mise en correspondance des émotions
  conversationnelles avec les états d'engagement, et la pondération des modalités
  dans la fusion — qui substitue une mesure d'impact à une simple déclaration de
  limite.

  \item Un \textbf{prototype logiciel complet et reproductible} (API, interface,
  scripts d'évaluation, journalisation des décisions), pensé comme un support de
  démonstration et d'expérimentation.

  \item Un \textbf{protocole d'évaluation multi-niveaux} documenté, prêt à être
  exécuté, articulant validation experte du modèle de décision, évaluation de la
  qualité du contenu et étude pilote auprès d'utilisateurs.
\end{enumerate}

%==============================================================================
\section{Méthodologie générale}
\label{sec:methodologie}

La démarche adoptée relève de la \emph{recherche orientée conception}
(\emph{design science research}), paradigme particulièrement adapté à la
conception d'artefacts informatiques répondant à un problème identifié
\cite{hevner2004design}. Elle articule l'analyse d'un problème réel — le
décrochage lié à la baisse d'engagement dans les MOOC —, la construction d'un
artefact — le framework \emolearn{} — et son évaluation au regard d'objectifs
explicites.

Concrètement, le travail s'organise autour d'une chaîne de traitement en trois
étages, illustrée par la figure~\ref{fig:pipeline} : un étage de
\emph{perception}, un étage de \emph{décision} et un étage de \emph{génération}.

\begin{figure}[H]
\centering
\begin{tikzpicture}[
    node distance=0.9cm and 1.2cm,
    box/.style={rectangle, rounded corners, draw=black, very thick,
                text width=3.4cm, minimum height=1.5cm, align=center,
                font=\small},
    io/.style={rectangle, draw=black!60, dashed, text width=3.4cm,
               minimum height=1.1cm, align=center, font=\footnotesize},
    arr/.style={-{Stealth[length=3mm]}, very thick}
]
\node[io] (in) {Scénario (clip vidéo + audio) + contexte pédagogique};
\node[box, below=of in] (perc) {\textbf{Perception}\\ Visage (DAiSEE) + Voix (IEMOCAP)\\ mappings appris + fusion};
\node[box, below=of perc] (dec) {\textbf{Décision}\\ état d'engagement\\ $\rightarrow$ type de contenu + patron};
\node[box, below=of dec] (gen) {\textbf{Génération}\\ LLM + 20 patrons de prompts};
\node[io, below=of gen] (out) {Contenu adaptatif + journal de décision};

\draw[arr] (in) -- (perc);
\draw[arr] (perc) -- (dec);
\draw[arr] (dec) -- (gen);
\draw[arr] (gen) -- (out);
\end{tikzpicture}
\caption{Vue d'ensemble de la chaîne de traitement d'\emolearn{} :
perception multimodale, décision pédagogique, puis génération de contenu
adaptatif.}
\label{fig:pipeline}
\end{figure}

\paragraph{Ressources de données.} Le projet s'appuie sur deux corpus publics de
référence, mobilisés à des fins complémentaires. Le corpus \textsc{DAiSEE}
\cite{gupta2016daisee} constitue le socle du projet : il fournit la taxonomie
d'engagement (quatre états déclinés en niveaux d'intensité) ainsi que les
scénarios d'apprentissage. Ce sont les \textbf{clips vidéo eux-mêmes} qui
alimentent la couche de perception, leurs annotations d'engagement étant réservées
au rôle de vérité terrain, pour l'entraînement des modèles supervisés et pour
l'évaluation. Le corpus \textsc{IEMOCAP} \cite{busso2008iemocap}
apporte la modalité vocale et, à travers un sous-corpus audiovisuel aligné, le
support de la validation de la fusion, seul cadre où le visage et la voix
proviennent effectivement du même locuteur.

\paragraph{Outils et environnement.} Le prototype repose sur une architecture
client-serveur : une interface développée avec une bibliothèque web moderne et un
service applicatif exposant une interface de programmation. La couche de
perception combine des modèles pré-entraînés — reconnaissance d'expressions
faciales et représentation neuronale de la parole — et des modèles supervisés
légers entraînés dans le cadre de ce travail, au moyen des bibliothèques usuelles
d'apprentissage automatique et d'apprentissage profond. La génération de
contenu s'appuie sur un grand modèle de langue accessible par interface de
programmation. L'ensemble est instrumenté de manière à journaliser chaque
décision (état détecté, stratégie retenue, patron mobilisé), garantissant la
traçabilité indispensable à l'analyse scientifique.

\paragraph{Stratégie d'évaluation.} Conformément à la nature exploratoire du
travail, l'évaluation est conçue comme un dispositif \emph{multi-niveaux}. Elle
combine des métriques automatiques (exactitude de la perception, indicateurs de
forme du contenu généré), une validation experte de la pertinence du modèle de
décision et du contenu produit, et une étude pilote auprès d'utilisateurs visant
à apprécier l'utilité et l'adéquation perçues. Ce protocole, détaillé au
chapitre~5, vise moins à démontrer une performance maximale qu'à rendre compte,
avec honnêteté, des forces et des limites du dispositif.

%==============================================================================
\section{Structure du mémoire}
\label{sec:structure}

Le présent mémoire s'organise en six chapitres.

Le \textbf{chapitre~1}, que conclut la présente section, a posé le contexte, la
problématique, les objectifs et la méthodologie générale du travail.

Le \textbf{chapitre~2} propose un état de l'art structuré autour de cinq axes :
l'engagement de l'apprenant dans les MOOC, la détection automatique de
l'engagement, la psychologie des émotions appliquée à l'apprentissage, l'IA
générative pour l'éducation et, enfin, les systèmes d'apprentissage adaptatifs.
Il se clôt sur une synthèse critique permettant de positionner \emolearn{} par
rapport aux travaux existants.

Le \textbf{chapitre~3} est consacré à la modélisation et à la conception du
framework. Il présente l'architecture globale, l'analyse du corpus \textsc{DAiSEE},
le modèle engagement\,$\rightarrow$\,adaptation — incluant l'apprentissage des
correspondances propres à chaque modalité et les deux familles de fusion —, le
système de génération fondé sur des patrons de prompts, ainsi que la conception du
prototype de simulation.

Le \textbf{chapitre~4} décrit l'implémentation : l'infrastructure technique, la
préparation des données, la réalisation du modèle de décision, du module de
génération adaptative et du prototype, jusqu'à la constitution du corpus de
contenu adaptatif.

Le \textbf{chapitre~5} expose la validation et les résultats. Il détaille la
méthodologie d'évaluation multi-niveaux, la validation du modèle de décision,
celle de la qualité du contenu généré, l'étude pilote utilisateur, puis discute
les résultats au regard de l'état de l'art.

Le \textbf{chapitre~6}, enfin, synthétise les contributions, répond aux questions
de recherche, expose les limites de l'étude et trace des perspectives, tant pour
la recherche que pour les implications pratiques dans le contexte des MOOC.

Le mémoire se termine par la bibliographie et un ensemble d'annexes regroupant
les patrons de prompts, les questionnaires, des exemples de contenu généré et les
résultats détaillés.
